\documentclass[11pt]{article}
\usepackage{amsmath,amssymb,amsthm,booktabs,hyperref}
\usepackage[margin=1in,marginparwidth=0.75in]{geometry}
\newtheorem{proposition}{Proposition}
\newtheorem{corollary}[proposition]{Corollary}
\newcommand{\old}{\marginpar{\scriptsize\textsf{Earlier theory}}}
\newcommand{\application}{\marginpar{\scriptsize\textsf{Application here}}}
\title{Finite orbits, small stabilizers, and generation}
\author{Human Compression: October corpus}
\date{}
\begin{document}
\maketitle

This note records the argument used for seven previously unresolved
generation tests in the fixed $G_2(11)$ computation. The orbit and
stabilizer argument is standard Schreier theory. The classification of
maximal subgroups and the description of the natural module are borrowed
explicitly below. The contribution here is their use with a bounded
certificate that checks whole matrices and distinguishes the two
generation outcomes by one membership test. No historical novelty is
claimed for the abstract group arguments.

The numerical assertions in the final table have separately executed
suppliers and strict consumers. Their independent execution audit and
canonical registration remain separate tasks. No rank-four intersection
property, whole manuscript recovery, open conjecture, or Lean
formalization is claimed by this note.

\section{An exact order certificate}

All matrices act on row vectors on the right. Let
$H=\langle g_0,g_1,g_2\rangle\leq\operatorname{GL}(V)$, where the $g_j$
are involutions, and let $v\ne0$. A rooted tree with $N$ vertices defines
words $T_0=1$ and $T_i=T_{p(i)}g_{j(i)}$, with $p(i)<i$.
Thus every $T_i$ is an actual word in $H$. Write $v_i=vT_i$.

\begin{proposition}\label{prop:certificate}\old
Suppose that the following conditions hold.
\begin{enumerate}
\item The vectors $v_0,\ldots,v_{N-1}$ are distinct.
\item A list $S=\{s_0=1,\ldots,s_{M-1}\}$ consists of distinct matrices
      fixing $v$. Every matrix in the list is an actual word in $H$,
      and $S$ is a subgroup of $H$.
\item For every $i$ and every $j\in\{0,1,2\}$, indices
      $k(i,j),\ell(i,j)$ satisfy the whole matrix equation
      \[
           T_i g_j=s_{\ell(i,j)}T_{k(i,j)}.
      \]
\end{enumerate}
Then $v^H=\{v_0,\ldots,v_{N-1}\}$, $S=H_v$, and
\[
                   H=\coprod_{i=0}^{N-1}S T_i,
                   \qquad |H|=MN.
\]
For a supplied matrix $x$, if $vx$ is absent from the vector list then
$x\notin H$. If $vx=v_k$, then $x\in H$ precisely when
$x=sT_k$ for some $s\in S$.
\end{proposition}

\begin{proof}
Put $\mathcal U=\bigcup_i S T_i$. This set contains the identity.
For $s\in S$, the displayed equation gives
\[
        (sT_i)g_j=s\,s_{\ell(i,j)}T_{k(i,j)}\in\mathcal U.
\]
Thus $\mathcal U$ contains every positive word in the involutions,
giving $H\subseteq\mathcal U$. Conversely $S\subseteq H$ and $T_i\in H$,
so $\mathcal U\subseteq H$.

If $sT_i=tT_k$, applying both matrices to $v$ gives $v_i=v_k$.
Distinctness implies $i=k$, and cancellation then gives $s=t$.
The union is disjoint and has $MN$ elements. Applying $v$ identifies its
complete orbit with the listed vectors. An element $sT_i$ fixes $v$
exactly when $i=0$, proving $H_v=S$. The same decomposition proves
both branches of the membership test.
\end{proof}

The subgroup condition on $S$ has its own finite certificate.
Choose at most eight genuine Schreier words
$h=T_i g_jT_{k(i,j)}^{-1}$. An earlier-parent tree proves that every
listed element of $S$ is a positive word in these $h$. A complete table
checks $sh\in S$ for every listed $s$ and every chosen $h$.
A finite positive monoid of invertible matrices is a group, so the list
is exactly their generated subgroup. Every listed matrix is checked
to fix $v$. This establishes condition~2 without assuming that the
seed stabilizer is cyclic or has an anticipated order.

Preservation of the vector orbit alone does not decide membership.
The final test compares whole matrices in the unique possible coset.
The factor $s$ in the edge equation is on the left; replacing $sT_k$
by $T_ks$ would not prove the closure used above.

\section{A maximality criterion from a prime element}

\begin{proposition}\label{prop:maximal}
Let $G\leq\operatorname{GL}_d(\mathbb F_q)$ be finite. Let $p$ be a
prime not dividing $q-1$, and let $m<|G|$. Suppose that every maximal
proper subgroup of $G$ whose order is divisible by $p$ either preserves
a one-dimensional $\mathbb F_q$ subspace or has order $m$.
Suppose $H\leq G$ satisfies
\[
 |H|=m,\qquad Y\in H,\qquad |Y|=p,\qquad
 \dim_{\mathbb F_q}\operatorname{Fix}(Y)=1.
\]
For a nonzero $v\in\operatorname{Fix}(Y)$, assume $|v^H|>q-1$.
Then $H$ is maximal in $G$. Consequently, for every $x\in G$,
\[
 \langle H,x\rangle=
 \begin{cases}
 H,&x\in H,\\
 G,&x\notin H.
 \end{cases}
\]
\end{proposition}

\begin{proof}
As $|H|<|G|$, choose a maximal proper subgroup $M$ containing $H$.
Its order is divisible by $p$. If $M$ preserves a rational line $L$,
then $Y$ acts on $L$ by a scalar $\lambda\in\mathbb F_q^\times$.
The order of $\lambda$ divides both $p$ and $q-1$, so $\lambda=1$.
Thus $L\subseteq\operatorname{Fix}(Y)$. Both spaces have dimension one,
whence $L=\operatorname{Fix}(Y)=\langle v\rangle$.
Then $v^H$ is contained in the $q-1$ nonzero vectors of $L$, a
contradiction. Therefore $|M|=m=|H|$, giving $H=M$.
The last assertion follows from maximality.
\end{proof}

This proposition combines a line argument and a maximal-overgroup
list. Its ambient-group assumptions require separate proofs; they do
not follow from a matrix dimension or a group name.

\section{The natural group \(G_2(11)\)}

\old
Retain the original split Cayley algebra over $\mathbb F_{11}$ and its
faithful action on the seven-dimensional pure subspace. Thus the ambient
group is the same $G=\operatorname{Aut}_{\mathbb F_{11}}(C)=G_2(11)$
used by the original computation. These model and membership statements
are explicit retained premises, not consequences of the order certificate.

The odd-characteristic maximal-subgroup classification of Kleidman
\cite{Kleidman}, in the full list restated before the exclusions in
Nguyen's proof of Theorem~4.4 \cite{Nguyen}, gives the following at
$q=11$: every proper maximal subgroup containing an element of
order~19 has type $\operatorname{SL}_3(11):2$ or $J_1$.
The other types have no factor~19 in their orders; there is no proper
subfield because 11 is prime. We use the full pre-exclusion list,
not the smaller conclusion of Nguyen's representation theorem.

The natural-module description of $\operatorname{SL}_3(11):2$ is
the stabilizer of a pure nonisotropic line; see Wilson, Section~4.2.6
\cite{Wilson}. The full outer extension preserves that line, although
it need not fix a chosen nonzero vector.
The other maximal type has order $|J_1|=175560$ \cite{Atlas}.
The standard order formula gives
\[
 |G_2(11)|=11^6(11^6-1)(11^2-1)>11^6=1771561>175560;
\]
see also \cite{BPP}.
Proposition~\ref{prop:maximal} applies with $p=19$, $m=175560$,
and $q-1=10$.

\begin{corollary}\label{cor:application}\application
Let $(\rho_0,\rho_1,\rho_2,\rho_3)$ be one of the seven original
four-tuples in Table~\ref{tab:seven}, in this same faithful natural model.
Put $H=\langle\rho_0,\rho_1,\rho_2\rangle$.
Under the retained maximal-subgroup and natural-module premises,
$H$ is maximal in $G$. The four tuples in the first line generate a
proper subgroup of $G$; the three in the second line generate $G$.
\end{corollary}

\begin{proof}
For each tuple a supplier and a separately executed strict consumer
establish Proposition~\ref{prop:certificate} with $N=3080$ and $M=57$.
Consequently $|H|=3080\cdot57=175560$.
Both roles freshly evaluate the specified original left-end word and
cofactor, proving $Y\ne1$, $Y^{19}=1$, and a one-dimensional fixed space.
The nonzero seed $v$ belongs to that space and has an orbit of
$3080>10$ distinct rows. Proposition~\ref{prop:maximal} makes $H$ maximal.
The strict whole-matrix membership test gives the values in the last
column of Table~\ref{tab:seven}. The full tuple generates
$\langle H,\rho_3\rangle$, so the two conclusions follow.
\end{proof}

\begin{table}[h]
\centering
\begin{tabular}{lll}
\toprule
Original positions in job $(20,4)$ & Strict order of $H$ & $\rho_3\in H$\\
\midrule
109, 349, 400, 502 & 175560 & yes\\
877, 1033, 1099    & 175560 & no\\
\bottomrule
\end{tabular}
\caption{Seven strict order and membership certificates.
Independent execution review and canonical adoption are separate.}
\label{tab:seven}
\end{table}

The three generating tuples still require a rank-four intersection
argument. Generation does not settle their original C-group
classification. The four nongenerating tuples are excluded from that
classification by generation, with no intersection claim.
The numerical workers leave their generation, named-type and
intersection-credit flags false; the corollary is the separate logical
application. The Schreier theorem, group orders and classification
remain attributed earlier work. No compression measurement or completed
original is added to the intake.

\begin{thebibliography}{9}
\bibitem{Kleidman}
P.~B.~Kleidman, The maximal subgroups of the Chevalley groups $G_2(q)$
with $q$ odd, the Ree groups \( {}^2G_2(q)\), and their automorphism groups,
\emph{J.~Algebra} \textbf{117} (1988), 30--71.
\url{https://doi.org/10.1016/0021-8693(88)90239-6}.
The classification is used through its retained restatement.
\bibitem{Nguyen}
H.~N.~Nguyen, \emph{Irreducible restrictions of Brauer characters of the
Chevalley group $G_2(q)$ to its proper subgroups},
arXiv:0910.4994v1, proof of Theorem~4.4, printed pp.~10--11.
\url{https://arxiv.org/pdf/0910.4994v1}.
\bibitem{Wilson}
R.~A.~Wilson, \emph{The Finite Simple Groups},
retained preliminary edition, Section~4.2.6, printed pp.~107--108.
\url{https://webspace.maths.qmul.ac.uk/r.a.wilson/TFSG.html}.
\bibitem{Atlas}
ATLAS authors, \emph{Janko group $J_1$}, group-order record.
\url{https://brauer.maths.qmul.ac.uk/Atlas/v3/spor/J1/}.
\bibitem{BPP}
J.~Bamberg, T.~Popiel and C.~E.~Praeger,
Simple groups, product actions, and generalized quadrangles,
\emph{Nagoya Math.~J.} \textbf{234} (2019), 87--126.
\url{https://doi.org/10.1017/nmj.2017.35};
preprint \url{https://arxiv.org/pdf/1702.07308}.
\end{thebibliography}
\end{document}

